\documentclass{article}

\usepackage[italian]{babel}
\usepackage[margin=20mm, footskip=20pt]{geometry}
\usepackage{graphicx}
\usepackage{subfiles}
\usepackage{hyperref}
\usepackage{nameref}
\usepackage{titlesec}
\usepackage{longtable}
\usepackage[table]{xcolor}
\usepackage{titling}
\usepackage{lastpage}
\usepackage{ifthen}
\usepackage{calc}
\usepackage{soulutf8}
\usepackage{contour}
\usepackage{float}
\usepackage{fancyhdr}
\usepackage{multirow}
\usepackage{pgfgantt}
\usepackage{lscape}
\usepackage{background}
\usepackage{lmodern}
\usepackage{textcomp}
\usepackage{lastpage}
\usepackage[utf8]{inputenc}
\usepackage{makecell}

% definizione dei percorsi in cui cercare immagini
\graphicspath{ {./}
    {./src/img/}
}

% setup della sottolineatura
\setuldepth{Flat}
\contourlength{0.8pt}

\newcommand{\uline}[1]{%
  \ul{{\phantom{#1}}}%
  \llap{\contour{white}{#1}}%
}


% setup dei link
\hypersetup{
  % set true if you want colored links (instead of boxes)
  colorlinks=true,
  % set to all if you want both sections and subsections linked
  linktoc=all,
  % set color for file links
  filecolor=blue,
  % set color for internal links
  linkcolor=black,
  % set url color
  urlcolor=blue,
  % set characters encoding in the bookmarks tab
  pdfencoding=unicode,
}

% setup forma \paragraph e \subparagraph
\titleformat{\paragraph}[hang]{\normalfont\normalsize\bfseries}{\theparagraph}{1em}{}
\titleformat{\subparagraph}[hang]{\normalfont\normalsize\bfseries}{\thesubparagraph}{1em}{}

% setup profondità indice di default
\setcounter{secnumdepth}{5}
\setcounter{tocdepth}{5}

\makeatletter %% non togliere, i comandi che definiscono i placeholder vanno qui
% esempio di utilizzo: \appendToGraphicspath{./img/} (un comando diverso per ogni path da includere)
% N.B.: ci DEVE essere un forward slash alla fine del path, a indicare che è una cartella.
\newcommand\appendToGraphicspath[1]{%
  \g@addto@macro\Ginput@path{{#1}}%
}

\newcommand{\setTitle}[1]{%
  \newcommand{\@placeholderTitle}{#1}%
}
\newcommand{\placeholderTitle}{\@placeholderTitle}

\newcommand{\setUso}[1]{%
  \newcommand{\@uso}{#1}%
}
\newcommand{\uso}{\@uso}

\newcommand{\setVersione}[1]{%
  \newcommand{\@versione}{#1}%
}
\newcommand{\versione}{\@versione}

\newcommand{\disabilitaVersione}{%
  \renewcommand{\setVersione}[1]{}%
  \renewcommand{\versione}{DISABILITATA}
}

\newcommand{\setResponsabile}[1]{%
  \newcommand{\@responsabile}{#1}%
}
\newcommand{\responsabile}{\@responsabile}

\newcommand{\setRedattori}[1]{%
  \newcommand{\@redattori}{#1}%
}
\newcommand{\redattori}{\@redattori}

\newcommand{\setVerificatori}[1]{%
  \newcommand{\@verificatori}{#1}%
}
\newcommand{\verificatori}{\@verificatori}

\newcommand{\setDestinatari}[1]{%
  \newcommand{\@destinatari}{#1}%
}
\newcommand{\destinatari}{\@destinatari}

\newcommand{\setDescrizione}[1]{%
  \newcommand{\@descrizione}{#1}%
}
\newcommand{\descrizione}{\@descrizione}

\newcommand{\setModifiche}[1]{%
  \newcommand{\@modifiche}{#1}%
}

\newcommand{\modifiche}{\@modifiche}
\makeatother %% non togliere, i comandi che definiscono i placeholder vanno qui

% hook per lo script che genera il glossario
\newcommand{\glossario}[1]{#1\textsubscript{G}}

% comandi per rendere opzionali gli elenchi di figure
\newcommand{\elencoFigure}{%
  \renewcommand{\listfigurename}{Elenco delle figure}%
  \listoffigures%
}

\newcommand{\disabilitaElencoFigure}{%
  \renewcommand{\elencoFigure}{}%
}

% comandi per rendere opzionali le tabelle
\newcommand{\elencoTabelle}{%
  \renewcommand{\listtablename}{Elenco delle tabelle}%
  \listoftables%
}

\newcommand{\disabilitaElencoTabelle}{%
  \renewcommand{\elencoTabelle}{}%
}

% comando per riferirsi ad una sezione
\newcommand{\refSec}[1]{%
  \hyperref[#1]{\S\ref*{#1}}%
}


\appendToGraphicspath{../../config/src/img/}

\setTitle{Verbale Esterno 2021-01-08}

\setVersione{v1.0.0}

\setResponsabile{
	Paparazzo Giorgia
}

\setRedattori{
	Rizzo Stefano
}

\setVerificatori{
	Kostadinov Samuel
}

\setUso{Esterno}

\setDestinatari{
	prof. Vardanega Tullio \\ &
	prof. Cardin Riccardo \\ &
	SpaghettiCode
}

\setDescrizione{Riassunto dell'incontro realizzato dal gruppo SpaghettiCode tenutosi l'8 gennaio 2021 in forma di meeting online con Piccoli Gregorio.}

\setModifiche{
	v1.0.0 & Paparazzo Giorgia & Responsabile & 2021-01-10 & Approvazione del documento & // \\
	v0.1.0 & Kostadinov Samuel & Verificatore & 2021-01-10 & Verifica del documento & // \\
	v0.0.1 & Rizzo Stefano & Analista & 2020-01-09 & Creazione del documento e prima stesura & //
}

\disabilitaElencoFigure
\disabilitaElencoTabelle

\begin{document}

\pagenumbering{gobble}


% variabile prima pagina (serve nell'istruzione dopo per stampare il bg solo nella prima pagina)
\newif\iffirstpage
\firstpagetrue

% setta l'immagina di bg
\backgroundsetup{
	scale=1,
	opacity=0.2,
	placement=top,
	contents={%
			\iffirstpage
				\includegraphics[width=\paperwidth]{datascience_og_colori.png}%
				\global\firstpagefalse
			\fi
		}%
}

\begin{titlepage}% per non stampare il numero della pagina

	\centering % allinea al centro la pagina
	\hspace{0.05\textwidth}% spazio tra linea e testo
	% lasciare questa riga per il corretto funzionamento di \parbox
	\parbox[b]{0.4\textwidth}{% cambiando la larghezza del testo il paragrafo si muove a destra o a sinistra 
	{\hspace{0.05\textwidth}\includegraphics[width=4cm,height=4cm]{logo_colori.png}}\\[3\baselineskip] % logo
	{\Huge\bfseries SpaghettiCode}\\ [\baselineskip] %titolo
	{\texttt{spaghetti.code.g6@gmail.com}}\\[\baselineskip]\\[4\baselineskip] % 
	{\Large\textsc\mbox{\placeholderTitle{}}}\\[4\baselineskip] % nome del documento
	{\begin{tabular}{r|l}
		\hline                                  \\
		% testo in grassetto
		\textbf{Versione}     & \versione{}     \\
		\rule{0pt}{3ex}%  EXTRA vertical height 
		\textbf{Approvazione} & \responsabile{} \\
		\rule{0pt}{3ex}%  EXTRA vertical height 
		\textbf{Redazione}    & \redattori{}    \\
		\rule{0pt}{3ex}%  EXTRA vertical height 
		\textbf{Verifica}     & \verificatori{} \\
		\rule{0pt}{3ex}%  EXTRA vertical height 
		\textbf{Uso}          & \uso{}          \\
		\rule{0pt}{3ex}%  EXTRA vertical height 
		\textbf{Destinato a}  & \destinatari{}
		\ifthenelse{\equal{\uso}{Esterno}}{
		\\ & Zucchetti S.p.A.
		}{}
	\end{tabular}}\\[4\baselineskip]

	}

	{\bfseries Descrizione}\\
	{\descrizione{}}\\[1\baselineskip]



\end{titlepage}

\newgeometry{textheight=660pt, lmargin=2cm, tmargin=2cm, rmargin=2cm}

% setup di header e footer nelle pagine senza numero
\fancypagestyle{nopage}{%
	\fancyhf{}%
	\fancyhead[R]{\includegraphics[width=1.3cm]{logo_colori.png}}%
	\fancyhead[L]{\emph{SpaghettiCode}\\\placeholderTitle{}}%
}
% setup di header e footer nelle pagine col numero
\fancypagestyle{usual}{%
	\fancyhf{}%
	\fancyhead[R]{\includegraphics[width=1.3cm]{logo_colori.png}}%
	\fancyhead[L]{\emph{SpaghettiCode}\\\placeholderTitle{}}%
	\fancyfoot[R]{\thepage\ di~\pageref{LastPage}}%
}
\setlength{\headheight}{1.8cm}

\newpage
\pagestyle{nopage}

\setcounter{table}{-1}

%REGISTRO DELLE MODIFICHE

\section*{Registro delle modifiche}
\label{sec:registro_delle_modifiche}

\rowcolors{2}{white!80!lightgray!90}{white}
\renewcommand{\arraystretch}{2} % allarga le righe con dello spazio sotto e sopra

\begin{longtable}
	[H]{|>{\centering\bfseries}m{2cm}|>{\centering}m{3.5cm}|>{\centering}m{2.5cm}|>{\centering}m{3cm}|>{\centering\arraybackslash}m{5cm}|}
	
	\hline
	\rowcolor{lightgray}
	{\textbf{Versione}} & {\textbf{Nominativo}} & {\textbf{Ruolo}} & {\textbf{Data}} & {\textbf{Descrizione}} \\
	\hline
	\endfirsthead
	
	\hline
	\rowcolor{lightgray}
	{\textbf{Versione}} & {\textbf{Nominativo}} & {\textbf{Ruolo}} & {\textbf{Data}} & {\textbf{Descrizione}} \\
	\hline
	\endhead
	
	\hline
	\rowcolor{white}
 	\multicolumn{5}{|c|}{\emph{Continua alla pagina successiva...}}\\
	\hline

	\endfoot
	\hline
	\endlastfoot
	
	\modifiche{}

\end{longtable}
% section registro_delle_modifiche (end)

\newpage
\thispagestyle{nopage}
\pagenumbering{roman}
\tableofcontents

\elencoFigure{}%
\elencoTabelle{}%

\newpage
\pagestyle{usual}
\pagenumbering{arabic}


\section{Informazioni generali}
\label{sec:info_generali}

\subsection{Informazioni incontro}
\label{sub:info_incontro}

\begin{itemize}
	\item \textbf{Luogo}: Applicazione desktop \glossario{Zoom};
	\item \textbf{Data}: 2021-01-09;
	\item \textbf{Ora}: 15:30-16:30
	\item \textbf{Partecipanti}:
	\begin{itemize}
		\item Contro Daniel Eduardo
		\item Fichera Jacopo
		\item Kostadinov Samuel
		\item Masevski Martin
		\item Pagotto Manuel
		\item Paparazzo Giorgia
		\item Rizzo Stefano
		\item Piccoli Gregorio
	\end{itemize}
\end{itemize}

\section{Ordine del giorno}%
\label{sec:ordine_del_giorno}

Di seguito vengono riportati i punti dell'ordine del giorno discussi:
\begin{itemize}
    \item Discussione funzionalità HD Viz;
    \item Chiarimenti presenza Database.
\end{itemize}


\subsection{Riferimenti}%
\label{sub:riferimenti}
\begin{itemize}
    \item Analisi dei requisiti versione 0.4.2;
	\item VE\_2020\_12\_17 1.0.0.
\end{itemize}

\section{Resoconto}
\label{sec:resoconto}

	L'incontro con il proponente è stato fissato al termine della prima stesura dell'Analisi dei Requisiti per discutere delle funzionalità  rilevate durante l'attività di analisi.
	Inoltre, l'incontro è servito anche a chiarire dei dubbi affiorati in seguito alla prima riunione con il proponente il 17/12/2020.
	Dopo una breve descrizione del funzionamento generale dell'applicativo, si ha proseguito l'incontro esponendo al proponente le varie funzionalità, discutendone l'utilità e eventuali miglioramenti.
	Di seguito sono elencati i punti più rilevanti trattati durante l'incontro.
	\begin{itemize}
	    \item È stato confermato che i dati da prelevare saranno contenuti in un database esterno. Questo conferma quanto detto nel capitolato e rettifica la riunione del 17/12.
		\item La funzionalità di modifica delle proprietà grafiche dei punti in uno Scatterplot matrix è stata rivista. Invece di evidenziare le dimensioni visualizzate con colori diversi, è stato deciso di utilizzare il colore come modo per rappresentare ulteriori dimensioni. A tale scopo sono stati introdotti gli use case:
		\begin {itemize}
			\item \textbf{UC4.2.2}: rappresentazione dimensione mediante tinta;
			\item \textbf{UC4.2.3}: rappresentazione dimensione mediante brillanza;
		\end{itemize}
		\item Nella visualizzazione Force Field è stato consigliato di aggiungere alcune funzionalità, tra cui:
		\begin {itemize}
		 	\item \textbf{preprocessing dei dati}: prima di effettuare il calcolo della matrice delle distanze, si potrà scegliere se normalizzare o standardizzare le dimensioni. In questa maniera sarà possibile evidenziare più facilmente diversi tipi di relazione tra i nodi.
			\item \textbf{assegnare pesi alle dimensioni}: poichè normalizzare o standardizzare le dimensioni può "appiattire" determinate relazioni tra dati, è stato consigliato di inserire un sistema per poter manualmente assegnare dei pesi alle dimensioni.
			\item \textbf{taglio di archi}: in diverse situazioni può essere vantaggioso eliminare collegamenti tra nodi che hanno forze di attrazione al di fuori di un certo intervallo.
		\end {itemize}
		\item La visualizzazione Heatmap è quella che ha subito più modifiche, questo perchè quella che il software "Orange canvas" chiama Heatmap, e che è stata usata come punto di partenza per l'analisi, non è la visualizzazione inizialmente richiesta dal proponente. Questa modalità di visualizzazione è stata modificata nel seguente modo:
		\begin {itemize}
			\item Quella che il proponente chiama "Heatmap" verrà obbligatoriamente implementata. Questa visualizzazione avrà il medesimo funzionamento della visualizzazione che il programma "Orange canvas" chiama "Distance map".
			\item La visualizzazione "Heatmap", com'era stata inizialmente intesa, è stata considerata interessante dal proponente. La sua implementazione è stata giudicata desiderabile.
			\item Entrambe le visualizzazioni dovranno prevedere una serie di funzionalità base come cambiare la scala di colori, ordinare righe e colonne mediante clustering gerarchico e visualizzazione delle label dei dati.
			\item Poichè le visualizzazioni Distance map e Force field prevedono entrambe l'utilizzo di una matrice delle distanze, è stato scelto di permettere le operazioni su suddetta matrice (preprocessing e assegnazione di pesi alle dimensioni) in entrambe le visualizzazioni.
		\end {itemize}
		\item Nella Proiezione Lineare Multi asse, viene espresso come "desiderabile" la funzionalità di poter interagire in tre dimensioni con il grafico.m
	  \end{itemize}

\section{Conclusione dell'incontro}
\label{sec:conclusione}
L'incontro ha chiarito i dubbi sulla presenza del database e ci ha permesso di correggere ed aggiungere gli use case nel documento Analisi dei Requisiti 0.4.2

\end{document}
